\section{Discussion}\label{sec:discussion}

\subsection{Why inverse temperatures, prices, and regularizers are the same kind of object}

Inverse temperature, market price, regularization strength, and attention-like precision are different things. In form, however, each is a multiplier that prices a limited resource in a layer that closes by constrained optimization \cite{jaynes1957information,arrowdebreu1954,sims2003rationalinattention}. Theorem~\ref{thm:duality} makes the shared form exact in the finite entropy setting. Whatever the substrate, once a ledger is budgeted and the layer selects the least-committed law within the budget, the Gibbs multiplier is a supporting price of that budget. When some weighted row has nonconstant costs, it is positive exactly when the budget lies strictly between the cheapest feasible cost and the cost of the unconstrained closure (Proposition~\ref{prop:regimes}).

This is more specific than saying that Lagrange multipliers exist. The usual statement starts from a given programme. The currency--constraint principle states where the programme's constraint comes from: a maintenance ledger together with a resource law (Assumption~\ref{ass:ledger}). It also states what happens to the price as the law tightens or relaxes. The principle is directional. Lower-layer spending becomes a higher-layer constraint, and the constraint becomes a higher-layer price. Layers are linked by maintenance, not merely by a formal duality.

\subsection{What the principle does not say}

The principle does not derive budgets from persistence. Example~\ref{ex:identity} shows that a perfectly persistent object can cost nothing, and an external ledger attached to it can be unbounded. A resource law is a separate physical, economic, or computational fact about the substrate. Conservation, a finite supply, a viability requirement, or a sustainability constraint can each provide one. Where such a law exists and the layer closes by maximum entropy with fixed row weights, a feasible budget has a price determined by Theorem~\ref{thm:duality} and Proposition~\ref{prop:regimes}. If some weighted row has nonconstant costs, the price is finite and nonnegative above the cheapest feasible cost, and no finite price attains that cost itself; otherwise the price is not identifiable. Where no resource law exists, the principle is silent.

Nor does the principle say that every layer closes by maximum entropy, or that one scalar price suffices. The laboratory uses a single-constraint Gibbs family because it is the simplest closure in which the price can be computed and tested. Several budgets would give a vector of prices, and other closures would give other duals.

\subsection{Capacity, drive, and resolution}

The ladder experiment makes a distinction that is easy to blur. A finer lens can expose more independent loops, and hence more coordinates in which a drive currency \emph{could} be expressed, without any additional drive being present. In the ladder construction the half-ring lens has five independent loops and zero affinity, because of a symmetry of the construction. The fine lens has eleven loops but only six driven directions, because the model has only six drive parameters. Coarsening can even create loops, as the path-to-triangle example shows. Statements about the ``dimension of the currency space'' at a given resolution should therefore say whether they mean available cycle coordinates or the rank of the actual drive.

\subsection{What counts as evidence here}

The laboratory is synthetic, and its role differs from claim to claim. For data processing, the price curve, and the packaging bound, the theorems carry the weight, and the numbers check that the implementation agrees with them. For the ladder, the ranks are exact properties of a specific construction, and the counterexample marks their limits. For the proxy comparison, the population statement (a fixed wrong kernel has positive expected excess loss) is a theorem, and the experiment shows a large gap at a modest sample size. The variability of fitted prices across samples is not a meaningful comparison between different cost coordinates, so it is not used as evidence for either cost.

Several modelling choices limit interpretation. The macro ledger uses the positive part of single-edge log-ratios, which is not invariant under a change of state potentials and is not the stationary entropy production. The macro Gibbs closures have complete support and need not correspond to transitions the substrate can realize. In the packaging experiment the price is chosen by an auxiliary calibration rather than inferred from the substrate's spending, and the substrate's actual spending is about \(50\) to \(300\) times lower. Uncertainty statements from the conditional-logit curvature are curvature scales, not sampling guarantees, under misspecification or at the boundary \(\lambda=0\).

\subsection{What this paper adds to the Six Birds series}

The foundations paper set out the vocabulary of closure: lenses, packaging, definability, and audits \cite{tsiokos2026sixbirds}. PICA supplied an interaction algebra and a laboratory in which primitives can be enabled and coupled \cite{tsiokos2026create}. This paper adds the spending side. It identifies what a closure must budget to keep its objects, gives the exact conditions under which that budget becomes a constraint and the constraint a price, and supplies audits that coarse observation cannot fake. The Lean development makes the audit inequality and the single-row duality machine-checked, and the packaging bound connects the price of leaving a block to how stable the block is as an object.
